% GENERATED FILE. Edit canonical lessons, then run npm run generate:books.
% canonical-source-hash: fnv1a64:a40c33557583869e
% canonical-lessons: FR-C62-revenir, FR-C62-etudier, FR-C62-payer, FR-C62-utiliser, FR-C62-montrer

\chapter{Come back, Study, Pay, Use, Show}
\label{ch:fr-come-back-study-pay-use-show}
\hlchaptermodality{62}

\begin{chapteropening}
\textbf{By the end of this chapter:} \emph{I can say to come back, to study, to pay, to use, to show.}

Complete the last lesson of chapter 62: I can say to come back, to study, to pay, to use, to show.
\end{chapteropening}

\section[revenir]{revenir — to come back}

\label{lesson:FR-C62-revenir}



\begin{quote}
Before the new one: say the French for to give, then the French for to live.
\end{quote}

\subsection*{You'll want to know: revenir}
\textbf{revenir} — \textquotedblleft{}to come back\textquotedblright{}.

\begin{cousinweb}
\textbf{revenir} $\leftarrow$ Latin \emph{revenire}, \textquotedblleft{}to come back\textquotedblright{}, from \emph{re-}, \textquotedblleft{}back\textquotedblright{}, and \emph{venire}, \textquotedblleft{}to come\textquotedblright{} — inherited. (EtymDB 2.1)
\end{cousinweb}

\begin{grammarlens}[title={What you've built}]
One more word for this chapter.
\end{grammarlens}

\subsection*{Your turn}
\begin{itemize}
\raggedright
  \item \emph{Say it:} \emph{revenir}
  \item \emph{Say it:} \emph{revenir}, once more
  \item \emph{Say it:} say \emph{revenir}
  \item \emph{Recall:} say the French for to give, then the French for to live, then say \emph{revenir} again
\end{itemize}

\begin{tcolorbox}[breakable,colback=teal!4,colframe=teal!35!black,title={Before you move on}]
What does revenir mean? (\textquotedblleft{}To come back\textquotedblright{}.) Say it once more.
\end{tcolorbox}
\section[étudier]{étudier — to study}

\label{lesson:FR-C62-etudier}



\begin{quote}
Before the new one: say the French for to live, then the French for to come back.
\end{quote}

\subsection*{You'll want to know: étudier}
\textbf{étudier} — \textquotedblleft{}to study\textquotedblright{}.

\begin{cousinweb}
\textbf{étudier} $\leftarrow$ Old French \emph{estudier}, a medieval scholars' word built on Latin \emph{studium}, \textquotedblleft{}eagerness, study\textquotedblright{}, from \emph{studere}, \textquotedblleft{}to be eager\textquotedblright{} — a learned borrowing. (EtymDB 2.1 records the Latin verb.) The \emph{é-} stands where an \emph{s} once was.
\end{cousinweb}

\begin{grammarlens}[title={What you've built}]
One more word for this chapter.
\end{grammarlens}

\subsection*{Your turn}
\begin{itemize}
\raggedright
  \item \emph{Say it:} \emph{étudier}
  \item \emph{Say it:} \emph{étudier}, once more
  \item \emph{Say it:} say \emph{étudier}
  \item \emph{Recall:} say the French for to live, then the French for to come back, then say \emph{étudier} again
\end{itemize}

\begin{tcolorbox}[breakable,colback=teal!4,colframe=teal!35!black,title={Before you move on}]
What does étudier mean? (\textquotedblleft{}To study\textquotedblright{}.) Say it once more.
\end{tcolorbox}
\section[payer]{payer — to pay}

\label{lesson:FR-C62-payer}



\begin{quote}
Before the new one: say the French for to come back, then the French for to study.
\end{quote}

\subsection*{You'll want to know: payer}
\textbf{payer} — \textquotedblleft{}to pay\textquotedblright{}.

\begin{cousinweb}
\textbf{payer} $\leftarrow$ Old French \emph{paiier} $\leftarrow$ Latin \emph{pacare}, \textquotedblleft{}to calm, to satisfy\textquotedblright{}, from \emph{pax}, \textquotedblleft{}peace\textquotedblright{} (EtymDB 2.1) — inherited: to pay was to settle with someone.
\end{cousinweb}

\begin{grammarlens}[title={What you've built}]
One more word for this chapter.
\end{grammarlens}

\subsection*{Your turn}
\begin{itemize}
\raggedright
  \item \emph{Say it:} \emph{payer}
  \item \emph{Say it:} \emph{payer}, once more
  \item \emph{Say it:} say \emph{payer}
  \item \emph{Recall:} say the French for to come back, then the French for to study, then say \emph{payer} again
\end{itemize}

\begin{tcolorbox}[breakable,colback=teal!4,colframe=teal!35!black,title={Before you move on}]
What does payer mean? (\textquotedblleft{}To pay\textquotedblright{}.) Say it once more.
\end{tcolorbox}
\section[utiliser]{utiliser — to use}

\label{lesson:FR-C62-utiliser}



\begin{quote}
Before the new one: say the French for to study, then the French for to pay.
\end{quote}

\subsection*{You'll want to know: utiliser}
\textbf{utiliser} — \textquotedblleft{}to use\textquotedblright{}.

\begin{cousinweb}
\textbf{utiliser} — formed in French in the 18th century from the word for \textquotedblleft{}useful\textquotedblright{}, a borrowing of Latin \emph{utilis}, \textquotedblleft{}useful\textquotedblright{}, from \emph{uti}, \textquotedblleft{}to use\textquotedblright{} (Lewis \& Short).
\end{cousinweb}

\begin{grammarlens}[title={What you've built}]
One more word for this chapter.
\end{grammarlens}

\subsection*{Your turn}
\begin{itemize}
\raggedright
  \item \emph{Say it:} \emph{utiliser}
  \item \emph{Say it:} \emph{utiliser}, once more
  \item \emph{Say it:} say \emph{utiliser}
  \item \emph{Recall:} say the French for to study, then the French for to pay, then say \emph{utiliser} again
\end{itemize}

\begin{tcolorbox}[breakable,colback=teal!4,colframe=teal!35!black,title={Before you move on}]
What does utiliser mean? (\textquotedblleft{}To use\textquotedblright{}.) Say it once more.
\end{tcolorbox}
\section[montrer]{montrer — to show}

\label{lesson:FR-C62-montrer}



\begin{quote}
Before the new one: say the French for to pay, then the French for to use.
\end{quote}

\subsection*{You'll want to know: montrer}
\textbf{montrer} — \textquotedblleft{}to show\textquotedblright{}.

\begin{cousinweb}
\textbf{montrer} $\leftarrow$ Old French \emph{monstrer} $\leftarrow$ Latin \emph{monstrare}, \textquotedblleft{}to show\textquotedblright{}, from \emph{monere}, \textquotedblleft{}to warn, to remind\textquotedblright{} — inherited. (EtymDB 2.1)
\end{cousinweb}

\begin{grammarlens}[title={What you've built}]
One more word for this chapter.
\end{grammarlens}

\subsection*{Your turn}
\begin{itemize}
\raggedright
  \item \emph{Say it:} \emph{montrer}
  \item \emph{Say it:} \emph{montrer}, once more
  \item \emph{Say it:} say \emph{montrer}
  \item \emph{Recall:} say the French for to pay, then the French for to use, then say \emph{montrer} again
\end{itemize}

\begin{tcolorbox}[breakable,colback=teal!4,colframe=teal!35!black,title={Before you move on}]
What does montrer mean? (\textquotedblleft{}To show\textquotedblright{}.) Say it once more.
\end{tcolorbox}
